\begin{definition}[The gather instance as a composition]
  \label{def:aba-gather-sub}
  \lean{PLTS.ABA.Gather.ProcRec}\lean{PLTS.ABA.Gather.GaNetState}\lean{PLTS.ABA.Gather.GaEvt}\lean{PLTS.ABA.Gather.inPull}\lean{PLTS.ABA.Gather.bindPull}\lean{PLTS.ABA.Gather.specPull}\lean{PLTS.ABA.Gather.SubStateAt}\lean{PLTS.ABA.Gather.instAt}\lean{PLTS.ABA.Gather.liftedSpec}\leanok
  \uses{def:aba-gather-spec, def:aba-brb-impl, def:aba-networkstate, def:syncproduct, def:process-algebra, def:idle-family, def:relabel}
  One gather instance over a payload type $X$: the $n$ gather programs beside the gather network, in parallel with $2n$ reliable-broadcast instances, the instance's five events hidden and the result read back over its interface alphabet --- \leandecl{PLTS.ABA.Gather.instAt}, which takes the broadcast tier as an argument.
  A program holds its own record, the messages delivered to it and two stores, the value each input-broadcast instance and the payload each bind-broadcast instance has returned here, written on the return events and read by the guards that consume them; the gather network holds the per-sender sent sets, the corrupted set and the instance's core (D29).
  Each broadcast instance is read along a pullback that names it, \leandecl{PLTS.ABA.Gather.inPull} or \leandecl{PLTS.ABA.Gather.bindPull}, so a label carrying another instance's index leaves that instance standing still, and the specification is read along \leandecl{PLTS.ABA.Gather.specPull}, which sends the call's input-enabledness loop to the call.
\end{definition}
